\chapter{KIỂM THỬ VÀ ĐÁNH GIÁ KẾT QUẢ}

\section{Môi trường kiểm thử}

Toàn bộ quá trình kiểm thử được thực hiện trên một máy tính cá nhân với cấu hình như sau:

\begin{table}[H]
\centering
\caption{Cấu hình môi trường kiểm thử}
\label{tab:test_environment}
\begin{tabular}{|l|p{8cm}|}
\hline
\textbf{Thành phần} & \textbf{Chi tiết} \\ \hline
Hệ điều hành & Windows 11 + WSL2 (Ubuntu 22.04) \\ \hline
CPU & AMD Ryzen 5 / Intel Core i5 (6 cores) \\ \hline
RAM & 16 GB DDR4 \\ \hline
Ổ cứng & SSD NVMe 512 GB \\ \hline
Docker & Docker Desktop 4.x (WSL2 backend) \\ \hline
Node.js & v20.x LTS \\ \hline
PostgreSQL & 15.x (chạy trong Docker container) \\ \hline
AvalancheGo & v1.11.x (chạy qua Avalanche Network Runner) \\ \hline
Trình duyệt & Google Chrome 120+ \\ \hline
\end{tabular}
\end{table}

\section{Kiểm thử chức năng}

\subsection{Kịch bản kiểm thử xác thực}

\begin{table}[H]
\centering
\caption{Kết quả kiểm thử chức năng xác thực}
\label{tab:test_auth}
\begin{tabular}{|c|p{4cm}|p{4cm}|c|}
\hline
\textbf{TC} & \textbf{Kịch bản} & \textbf{Kết quả mong đợi} & \textbf{Kết quả} \\ \hline
TC-01 & Đăng ký tài khoản mới & Tạo thành công, trả về user info & \textcolor{green}{PASS} \\ \hline
TC-02 & Đăng ký email đã tồn tại & Trả về lỗi 409 Conflict & \textcolor{green}{PASS} \\ \hline
TC-03 & Đăng nhập đúng email/password & Trả về JWT token + user info & \textcolor{green}{PASS} \\ \hline
TC-04 & Đăng nhập sai password & Trả về lỗi 401 Unauthorized & \textcolor{green}{PASS} \\ \hline
TC-05 & Truy cập API không có JWT & Trả về lỗi 401 Unauthorized & \textcolor{green}{PASS} \\ \hline
TC-06 & Truy cập API với JWT hết hạn & Trả về lỗi 401 Unauthorized & \textcolor{green}{PASS} \\ \hline
TC-07 & Đăng nhập Google OAuth & Redirect Google $\rightarrow$ callback $\rightarrow$ tạo user + JWT & \textcolor{green}{PASS} \\ \hline
\end{tabular}
\end{table}

\subsection{Kịch bản kiểm thử quản lý Subnet}

\begin{table}[H]
\centering
\caption{Kết quả kiểm thử chức năng Subnet}
\label{tab:test_subnet}
\begin{tabular}{|c|p{4cm}|p{4cm}|c|}
\hline
\textbf{TC} & \textbf{Kịch bản} & \textbf{Kết quả mong đợi} & \textbf{Kết quả} \\ \hline
TC-08 & Tạo Subnet mới & Subnet được tạo, validator container khởi động & \textcolor{green}{PASS} \\ \hline
TC-09 & Xem danh sách Subnet & Hiển thị đúng Subnet thuộc user hiện tại & \textcolor{green}{PASS} \\ \hline
TC-10 & Dừng Subnet & Container validator dừng, trạng thái chuyển ``stopped'' & \textcolor{green}{PASS} \\ \hline
TC-11 & Khởi động lại Subnet & Container khởi động lại, RPC endpoint hoạt động & \textcolor{green}{PASS} \\ \hline
TC-12 & Xóa Subnet & Container bị xóa, bản ghi DB bị xóa & \textcolor{green}{PASS} \\ \hline
TC-13 & Xóa Subnet của người khác & Trả về lỗi 403 Forbidden & \textcolor{green}{PASS} \\ \hline
\end{tabular}
\end{table}

\subsection{Kịch bản kiểm thử giám sát}

\begin{table}[H]
\centering
\caption{Kết quả kiểm thử giám sát thời gian thực}
\label{tab:test_monitoring}
\begin{tabular}{|c|p{4cm}|p{4cm}|c|}
\hline
\textbf{TC} & \textbf{Kịch bản} & \textbf{Kết quả mong đợi} & \textbf{Kết quả} \\ \hline
TC-14 & Kết nối WebSocket & Gateway nhận connection, bắt đầu broadcast & \textcolor{green}{PASS} \\ \hline
TC-15 & Nhận dữ liệu status & Client nhận event ``status'' mỗi 5 giây & \textcolor{green}{PASS} \\ \hline
TC-16 & Subscribe chain cụ thể & Client nhận ``chain-stats'' mỗi 3 giây & \textcolor{green}{PASS} \\ \hline
TC-17 & Node offline & Dashboard hiển thị ``Node offline'' thay vì crash & \textcolor{green}{PASS} \\ \hline
TC-18 & Ngắt kết nối client & Interval được cleanup, không memory leak & \textcolor{green}{PASS} \\ \hline
\end{tabular}
\end{table}

\section{Đánh giá hiệu năng}

\subsection{Thời gian triển khai Subnet}

Em đo thời gian triển khai Subnet bằng hai phương thức: thủ công (CLI) và tự động (qua hệ thống), mỗi phương thức lặp lại 3 lần lấy trung bình.

\begin{table}[H]
\centering
\caption{So sánh thời gian triển khai Subnet: thủ công vs. tự động}
\label{tab:deployment_time}
\begin{tabular}{|l|c|c|c|c|}
\hline
\textbf{Phương thức} & \textbf{Lần 1} & \textbf{Lần 2} & \textbf{Lần 3} & \textbf{Trung bình} \\ \hline
Thủ công (CLI + cấu hình) & 45 phút & 38 phút & 42 phút & $\sim$42 phút \\ \hline
Tự động (Web Wizard) & 1p 48s & 2p 05s & 1p 58s & $\sim$1p 57s \\ \hline
\textbf{Tỷ lệ cải thiện} & \multicolumn{4}{c|}{\textbf{Giảm $\sim$95\% thời gian}} \\ \hline
\end{tabular}
\end{table}

Kết quả cho thấy hệ thống giảm đáng kể thời gian triển khai --- từ trung bình 42 phút (thao tác CLI, cấu hình file, đợi đồng bộ) xuống còn dưới 2 phút thông qua wizard giao diện web. Yếu tố chiếm thời gian chính trong phương thức tự động là thời gian Docker pull image và khởi động container.

\subsection{Hiệu năng API Response Time}

\begin{table}[H]
\centering
\caption{Thời gian phản hồi API trung bình (đo 50 request)}
\label{tab:api_performance}
\begin{tabular}{|l|c|c|c|}
\hline
\textbf{Endpoint} & \textbf{Avg (ms)} & \textbf{P95 (ms)} & \textbf{P99 (ms)} \\ \hline
\texttt{POST /auth/login} & 85 & 120 & 180 \\ \hline
\texttt{GET /subnets} & 25 & 45 & 65 \\ \hline
\texttt{GET /node-status/dashboard} & 150 & 280 & 450 \\ \hline
\texttt{GET /node-status/blocks} & 180 & 320 & 500 \\ \hline
\texttt{POST /subnets (create)} & 2.500 & 5.000 & 8.000 \\ \hline
\texttt{POST /contracts/deploy} & 3.200 & 6.500 & 12.000 \\ \hline
\end{tabular}
\end{table}

Hầu hết API phản hồi dưới 300ms (P95), đáp ứng yêu cầu NFR-01. Riêng hai endpoint tạo Subnet và deploy Contract có thời gian dài hơn do phải tương tác với Docker API và Blockchain node --- đây là các thao tác đặc thù không thể tối ưu thêm về mặt logic.

\subsection{Khả năng mở rộng --- Số Subnet tối đa}

Em tiến hành thử nghiệm tăng dần số Subnet chạy đồng thời trên máy cấu hình 16GB RAM để tìm điểm gãy (breaking point):

\begin{table}[H]
\centering
\caption{Khả năng mở rộng: Số Subnet đồng thời trên 16GB RAM}
\label{tab:scalability}
\begin{tabular}{|c|c|c|c|}
\hline
\textbf{Số Subnet} & \textbf{RAM sử dụng} & \textbf{CPU Load} & \textbf{Trạng thái} \\ \hline
1 & 2.1 GB & 15\% & Ổn định \\ \hline
3 & 5.8 GB & 35\% & Ổn định \\ \hline
5 & 9.2 GB & 55\% & Ổn định \\ \hline
7 & 12.5 GB & 72\% & Chậm nhẹ \\ \hline
10 & 15.8 GB & 95\% & Bắt đầu swap, rất chậm \\ \hline
\end{tabular}
\end{table}

Kết quả cho thấy với cấu hình 16GB RAM, hệ thống hoạt động ổn định với 5 Subnet đồng thời. Từ Subnet thứ 7 trở đi bắt đầu xuất hiện hiện tượng chậm do cạnh tranh tài nguyên. Điểm gãy xảy ra ở khoảng 10 Subnet khi hệ thống bắt đầu sử dụng swap memory.

\section{Đánh giá tổng thể}

\subsection{Ưu điểm}

\begin{enumerate}
    \item Giao diện web trực quan, dễ sử dụng, không yêu cầu kiến thức CLI chuyên sâu.
    \item Giảm 95\% thời gian triển khai Subnet so với phương thức thủ công.
    \item Giám sát real-time qua WebSocket với độ trễ dưới 3 giây.
    \item Xác thực và phân quyền chặt chẽ (JWT + Google OAuth + ownership check).
    \item Kiến trúc modular dễ mở rộng, thêm module mới không ảnh hưởng module hiện tại.
    \item Graceful degradation --- hệ thống hoạt động bình thường ngay cả khi Avalanche node offline.
\end{enumerate}

\subsection{Hạn chế}

\begin{enumerate}
    \item Chưa hỗ trợ triển khai trên nhiều máy chủ phân tán (multi-host deployment); hiện tại giới hạn trên một máy.
    \item Chưa có cơ chế auto-scaling tự động (tăng/giảm validator dựa trên tải mạng).
    \item Phụ thuộc vào Docker Desktop trên Windows/macOS; cần WSL2 hoặc Linux native cho hiệu năng tốt nhất.
    \item Chưa tích hợp hệ thống thông báo qua email/Telegram khi có cảnh báo nghiêm trọng.
    \item Giao diện quản lý Prometheus/Grafana vẫn cần truy cập riêng, chưa nhúng trực tiếp vào Dashboard.
\end{enumerate}
